% =====================================================================================
% REVISIÓN 1  -  Introducción, Planteamiento del problema, Objetivos, Justificación, Alcances y limitaciones
% Escribe tu texto FUERA de los recuadros \begin{guia} ... \end{guia}.
% No cambies los títulos ni borres los comentarios «% audit:req=...» al final de cada \section / \subsection.
% =====================================================================================

\section{Introducción} \label{sec:introduccion} % audit:req=introduccion

\begin{guia}[Guía: Introducción (Revisión 1)]
\textbf{¿Para qué sirve?} Presenta tu trabajo a un lector que no lo conoce y lo invita a seguir leyendo. Es el resumen
``de entrada'' del proyecto: contexto, problema, propuesta y organización del documento.

\textbf{Debe responder:}
\begin{itemize}[nosep,leftmargin=*]
  \item ¿En qué contexto (sector, empresa, área tecnológica) se desarrolla el proyecto?
  \item ¿Cuál es el problema o la necesidad, dicho en pocas líneas? (el detalle va en \emph{Planteamiento del problema}).
  \item ¿Qué se propone hacer y qué se espera lograr?
  \item ¿Por qué es importante? (una idea; el detalle va en \emph{Justificación}).
  \item ¿Cómo está organizado el documento? (una frase por sección o capítulo).
\end{itemize}

\textbf{Extensión mínima:} 150 palabras (1 a 2 páginas), redactadas con tus propias palabras.

\textbf{Errores frecuentes:} copiar definiciones de Internet; escribir solo ``en este capítulo se verá\ldots'' sin decir
nada concreto; dejar la introducción para el final y no actualizarla cuando el proyecto cambia.

\textbf{Cómo se revisa:} claridad del contexto, que el problema y la propuesta se entiendan sin leer el resto, y que la
organización descrita coincida con las secciones reales del documento.
\end{guia}

% >>> ESCRIBE AQUÍ el texto de la introducción (fuera de la guía) <<<

La administración pública estatal sostiene su funcionamiento sobre el documento. Cada trámite que un ciudadano inicia y cada resolución que una autoridad emite quedan asentados en un soporte documental que debe conservarse, localizarse y, llegado el caso, comprobarse ante instancias de control. Cuando ese soporte circula en papel y su control se lleva de manera manual, la institución conserva los documentos pero pierde el rastro de lo que ocurre con ellos.

El Instituto Registral y Catastral del Estado de Tamaulipas es un organismo público descentralizado de la administración pública estatal, con personalidad jurídica y patrimonio propio, que tiene por objeto desarrollar las tareas del Registro Público de la Propiedad Inmueble y del Comercio, así como del Catastro del Estado, mediante la coordinación entre ambos servicios \cite{poe2022irct}. Su función sustantiva consiste en otorgar certeza jurídica sobre los actos relacionados con el derecho de propiedad, lo que implica que la documentación que resguarda no es meramente administrativa: constituye el respaldo material de derechos patrimoniales de los particulares.

El Instituto recibe y resguarda un volumen constante de documentación proveniente de sus distintas oficinas. Ese flujo depende actualmente del traslado físico de los expedientes, de la entrega presencial y del registro manual de lo recibido, lo que provoca demoras en la atención, riesgo de extravío o duplicidad y, sobre todo, la imposibilidad de conocer con precisión cuánto se está procesando, quién lo procesa y desde qué oficina.

El presente trabajo describe la integración de un sistema que centraliza este proceso en un solo entorno digital. El usuario registra su solicitud y carga los documentos, los cuales quedan almacenados con sus metadatos y pueden consultarse en línea sin necesidad de descargarlos, bajo un esquema de permisos diferenciados según el perfil y la oficina de adscripción. Cada movimiento queda asentado en una bitácora que garantiza la trazabilidad y que alimenta un módulo de control estadístico capaz de mostrar el volumen documental por usuario, oficina, tipo de documento y periodo, con reportes exportables para la toma de decisiones.

La solución no se limita al desarrollo del software: contempla también la creación y configuración del entorno virtualizado que la aloja, la definición de políticas de seguridad, respaldo y monitoreo, y su liberación a producción. Esta dimensión de infraestructura sustenta el carácter de integración de tecnologías de información del proyecto, ya que el resultado entregado es un sistema operando en un ambiente real y administrado, y no un prototipo aislado.

El documento se organiza de la siguiente manera. La presente sección plantea el problema, fija los objetivos, justifica el trabajo y delimita su alcance. La Sección~\ref{sec:antecedentes} describe al organismo receptor y revisa las plataformas de gestión documental existentes frente a la propuesta. La Sección~\ref{sec:marco_teorico} reúne los conceptos y estándares necesarios para comprender el desarrollo. La Sección~\ref{sec:propuesta} expone la arquitectura y el modelo de datos de la solución, la Sección~\ref{sec:desarrollo_sistema} documenta su construcción y pruebas, y la Sección~\ref{sec:conclusiones} presenta las conclusiones y el trabajo futuro.



\subsection{Planteamiento del problema} \label{sec:problema} % audit:req=planteamiento

\begin{guia}[Guía: Planteamiento del problema (Revisión 1)]
\textbf{¿Para qué sirve?} Describe con precisión la situación que quieres resolver y demuestra que es un problema real.

\textbf{Debe responder:}
\begin{itemize}[nosep,leftmargin=*]
  \item ¿Qué situación actual es insatisfactoria o genera una necesidad? Descríbela con hechos, ejemplos o cifras.
  \item ¿Quién la padece (usuarios, área, empresa) y cómo se manifiesta (retrasos, errores, costos, quejas)?
  \item ¿Qué consecuencias tiene no resolverla?
  \item ¿Qué se ha intentado hasta ahora y qué falta (la brecha)? Conecta con los \emph{Antecedentes}.
  \item ¿Cuál es el problema en una sola frase? Escríbelo como enunciado del problema o como pregunta de
        investigación (p. ej. ``¿Cómo reducir X en Y mediante Z?'').
\end{itemize}

\textbf{Extensión mínima:} 120 palabras.

\textbf{Errores frecuentes:} describir la solución en lugar del problema (``el problema es que no hay un sistema web'');
problemas sin evidencia; enunciados tan amplios que no se pueden resolver en una estadía.

\textbf{Cómo se revisa:} que el problema esté delimitado, sea verificable y que los objetivos lo atiendan.
\end{guia}

% >>> ESCRIBE AQUÍ el planteamiento del problema <<<

El proceso documental vigente se apoya en tres mecanismos desarticulados: el traslado físico del expediente entre oficinas, la entrega presencial con sello y firma de recibido, y el registro manual en libretas u hojas de cálculo administradas localmente por cada oficina. Este esquema no es una elección deliberada, sino el resultado de que el flujo administrativo interno nunca fue formalizado en un procedimiento documentado, de modo que cada oficina resolvió su control de la manera que le resultó práctica.

De ahí se derivan cuatro cuellos de botella. El traslado físico somete el tiempo de ciclo del documento a la logística de transporte y no al trabajo sustantivo que debe realizarse sobre él. La doble captura obliga a registrar los mismos datos en la oficina emisora y en la receptora, lo que duplica el esfuerzo e introduce discrepancias entre ambos registros. La búsqueda en archivo depende del conocimiento tácito del encargado, de modo que su ausencia detiene la consulta. Y la integración manual de reportes produce información tardía que no puede verificarse contra la fuente que la originó.

Por encima de ellos está el problema central, que no es la ausencia de digitalización sino la ausencia de trazabilidad. Del análisis del flujo se desprenden cinco puntos donde la cadena de custodia de la información se interrumpe:

\begin{enumerate}[nosep]
  \item \textbf{Tránsito invisible.} Entre la salida de la oficina emisora y el sello de recepción no se registra el estado del documento, por lo que durante horas o días nadie puede afirmar dónde se encuentra el expediente.
  \item \textbf{Ausencia de identificador único.} Cada oficina folia conforme a su propio criterio, lo que impide seguir un mismo documento entre oficinas y favorece la duplicidad.
  \item \textbf{Registro heterogéneo.} Los campos capturados varían entre oficinas y entre personas, de modo que los datos no resultan comparables ni agregables.
  \item \textbf{Consulta sin huella.} El acceso físico al expediente no deja constancia de quién consultó ni con qué motivo, lo que impide auditar el acceso a documentación sensible.
  \item \textbf{Sustitución sin versionado.} El documento corregido reemplaza al anterior sin conservar el antecedente, y se pierde la posibilidad de comprobar qué se tenía en una fecha determinada.
\end{enumerate}

Conviene precisar el origen de esta caracterización. El proceso vigente no está formalizado en ningún procedimiento documentado del Instituto, por lo que su descripción se construyó a partir del análisis del flujo elaborado con el asesor empresarial durante la definición del proyecto. Los cinco puntos de ruptura son, en consecuencia, resultado de analizar ese flujo y no de una medición previa: el proceso actual no produce los datos que harían posible medirlo, lo cual es en sí mismo el problema que se plantea. Esa es precisamente la razón por la que el sistema incorpora una bitácora desde su primera versión, para que a partir de su liberación el Instituto disponga por primera vez de información verificable sobre su propio flujo documental.

La consecuencia conjunta es que el Instituto conserva sus documentos pero no puede dar cuenta de su ciclo de vida. Esto adquiere particular relevancia en una institución cuya documentación soporta derechos de propiedad y que, como sujeto obligado, está sometida a la normativa en materia de archivos \cite{lga2018} y de protección de datos personales \cite{lgpdppso2017}. De no atenderse, la institución continuará sin poder responder con certeza dónde está un expediente, quién accedió a él ni cuánto trabajo documental realiza cada oficina, y seguirá tomando decisiones de carga de trabajo sobre cifras que no son verificables.

El problema se enuncia así: \textbf{¿de qué manera puede integrarse un sistema que automatice el proceso de solicitud, carga, consulta y control estadístico de documentos en el Instituto Registral y Catastral del Estado de Tamaulipas, de modo que se reduzcan los tiempos de atención y se garantice la trazabilidad documental sobre una infraestructura administrable por la propia institución?}



\subsection{Objetivos} \label{sec:objetivos} % audit:req=objetivos

\begin{guia}[Guía: Objetivo general y objetivos específicos (Revisión 1)]
\textbf{¿Para qué sirve?} Fija qué vas a lograr. Todo lo demás (resultados y conclusiones) se compara contra estos objetivos.

\textbf{Objetivo general} (uno): verbo en infinitivo + qué se hará + con qué o cómo + para qué, y si es posible un criterio
medible. \emph{Ejemplo:} ``Desarrollar un sistema web de notificaciones para familiares de pacientes hospitalizados que
reduzca en 50\% las llamadas de consulta al área de enfermería.''

\textbf{Objetivos específicos} (de 3 a 6): pasos concretos y verificables que, al cumplirse, logran el objetivo general.
Cada uno debe poder demostrarse con un resultado (un módulo, un diagrama, un experimento) que luego aparecerá en
\emph{Resultados} y se retomará en \emph{Conclusiones}. Deben ser SMART: específicos, medibles, alcanzables, relevantes y
con tiempo definido.

\textbf{Debe responder:} ¿qué voy a entregar? ¿cómo sabré que lo logré? ¿cada objetivo específico aporta al general?

\textbf{Extensión mínima:} 40 palabras en total.

\textbf{Errores frecuentes:} verbos no medibles (``conocer'', ``entender'', ``estudiar''); objetivos que son actividades
(``investigar sobre\ldots'') y no logros; más de un objetivo general; objetivos que luego no se atienden en el documento.

\textbf{Cómo se revisa:} coherencia con el problema, medibilidad y que cada objetivo tenga un resultado asociado.
\end{guia}

\subsubsection{Objetivo general} \label{subsec:objetivo_general}
% >>> ESCRIBE AQUÍ el objetivo general (una sola oración) <<<

Integrar un sistema que automatice el proceso de solicitud, carga, consulta y control estadístico de documentos en el Instituto Registral y Catastral del Estado de Tamaulipas, mediante su desarrollo e implementación sobre un entorno virtualizado, con el fin de reducir los tiempos de atención, garantizar la trazabilidad documental y generar información confiable para la toma de decisiones.


\subsubsection{Objetivos específicos} \label{subsec:objetivos_especificos}
% >>> ESCRIBE AQUÍ los objetivos específicos como lista numerada: \begin{enumerate} \item ... \end{enumerate} <<<

\begin{enumerate}
  \item \textbf{Diagnosticar} el proceso vigente de solicitud y resguardo documental, identificando tiempos de atención, cuellos de botella y puntos donde se pierde la trazabilidad. \textit{Verificable con} el documento de diagnóstico y la caracterización del flujo.
  \item \textbf{Especificar} los requerimientos funcionales y no funcionales del sistema con los responsables de cada oficina, y definir con ellos el alcance, la metodología de desarrollo y el cronograma. \textit{Verificable con} el catálogo de requerimientos y las actas de validación firmadas.
  \item \textbf{Diseñar} la arquitectura de la solución, el modelo de datos que soporta el registro y la trazabilidad documental, y los prototipos de interfaz de los procesos de solicitud, carga, consulta y visualización de indicadores. \textit{Verificable con} el diagrama entidad-relación, el documento de arquitectura y los prototipos validados con usuarios.
  \item \textbf{Implementar} el entorno virtualizado que aloja la solución y el entorno de desarrollo, incluyendo sistema operativo, servicios base, red, usuarios, gestor de base de datos y control de versiones. \textit{Verificable con} el entorno operando y el repositorio con historial de cambios.
  \item \textbf{Desarrollar} los cuatro módulos funcionales del sistema: autenticación con permisos diferenciados por perfil y oficina de adscripción; solicitud y carga de documentos con validación y captura de metadatos; consulta y visualización en línea con bitácora de accesos; y control estadístico con reportes exportables. \textit{Verificable con} cada módulo en operación y las pruebas por historia de usuario.
  \item \textbf{Liberar y transferir} la solución, ejecutando las pruebas funcionales, de integración y de seguridad, poniéndola en producción con respaldos y monitoreo, y entregando los manuales técnico y de usuario junto con la capacitación al personal. \textit{Verificable con} las evidencias de prueba, el acta de liberación y la lista de asistencia a la capacitación.
\end{enumerate}

Los seis objetivos específicos agrupan, sin omitir ninguna, las catorce actividades comprometidas en el registro del proyecto ante la Universidad. La Tabla~\ref{tab:objetivos_actividades} establece esa correspondencia, de modo que cada actividad registrada tenga un objetivo que la respalde y ningún objetivo carezca de actividad que lo ejecute.

\begin{table}[H]
  \centering
  \footnotesize
  \setlength{\tabcolsep}{5pt}
  \renewcommand{\arraystretch}{1.2}
  \caption{Correspondencia entre los objetivos específicos y las actividades registradas del proyecto.}
  \label{tab:objetivos_actividades}
  \begin{tabular}{@{}cp{10.6cm}c@{}}
    \toprule
    \textbf{Objetivo} & \textbf{Actividad registrada} & \textbf{Núm.} \\
    \midrule
    1 & Diagnosticar el proceso actual de solicitud y resguardo documental, identificando tiempos de atención, cuellos de botella y puntos donde se pierde la trazabilidad & 1 \\
    2 & Levantar los requerimientos funcionales y no funcionales del sistema con los usuarios responsables de cada oficina & 2 \\
    2 & Definir el alcance, los objetivos específicos, la metodología de desarrollo y el cronograma de trabajo del proyecto & 3 \\
    3 & Diseñar la arquitectura de la solución y modelar la base de datos que soportará el registro y la trazabilidad documental & 4 \\
    3 & Diseñar los prototipos de interfaz para los procesos de solicitud, carga, consulta y visualización de indicadores & 5 \\
    4 & Crear y configurar el entorno virtualizado que alojará la solución, incluyendo sistema operativo, red, servicios base y usuarios & 6 \\
    4 & Configurar el entorno de desarrollo, el gestor de base de datos y el control de versiones del proyecto & 7 \\
    5 & Desarrollar el módulo de autenticación con roles y permisos diferenciados según el perfil y la oficina de adscripción del usuario & 8 \\
    5 & Desarrollar el módulo de solicitud y carga de documentos con validación de formato, tamaño y captura de metadatos & 9 \\
    5 & Implementar el componente de consulta y visualización de documentos en línea con registro de bitácora de accesos & 10 \\
    5 & Desarrollar el módulo de control estadístico que muestre el volumen documental por usuario, oficina, tipo de documento y periodo, con reportes exportables & 11 \\
    6 & Ejecutar las pruebas funcionales, de integración y de seguridad, corrigiendo las incidencias detectadas & 12 \\
    6 & Liberar el sistema a producción en el entorno virtualizado, configurando respaldos, monitoreo y realizando pruebas de operación real & 13 \\
    6 & Elaborar los manuales técnico y de usuario, capacitar al personal e integrar el reporte final de estadías & 14 \\
    \bottomrule
  \end{tabular}
\end{table}

\subsection{Justificación} \label{sec:justificacion} % audit:req=justificacion

\begin{guia}[Guía: Justificación (Revisión 1)]
\textbf{¿Para qué sirve?} Explica por qué vale la pena dedicar tiempo y recursos a este proyecto.

\textbf{Debe responder:}
\begin{itemize}[nosep,leftmargin=*]
  \item ¿Quién se beneficia y de qué manera? (empresa, usuarios, la comunidad, tú).
  \item ¿Qué se gana en términos prácticos, económicos, sociales, tecnológicos o científicos? Usa datos cuando existan.
  \item ¿Qué pasaría si no se hace?
  \item ¿Por qué es viable hacerlo ahora y por qué con este enfoque?
  \item ¿Cómo se relaciona con el perfil de egreso de tu programa académico?
\end{itemize}

\textbf{Extensión mínima:} 100 palabras.

\textbf{Errores frecuentes:} repetir el problema con otras palabras; afirmaciones sin sustento (``mejorará mucho la
eficiencia''); olvidar al beneficiario.

\textbf{Cómo se revisa:} que los beneficios sean concretos y consistentes con el problema y los objetivos.
\end{guia}

% >>> ESCRIBE AQUÍ la justificación <<<

El Instituto ejecuta actualmente un Proyecto Ejecutivo de Modernización Integral orientado a optimizar sus procesos técnicos, administrativos y tecnológicos \cite{tamaulipas2025mod}. El presente proyecto se inserta en esa línea de trabajo atendiendo el flujo administrativo interno de la documentación, que es la capa sobre la que descansan los procesos sustantivos ya modernizados. Conviene precisar el alcance de esta aportación: el Registro Público del Estado opera desde 2019 bajo un sistema registral electrónico que atiende los trámites registrales de forma digital \cite{futuver2024}. Lo que este proyecto atiende no es ese proceso sustantivo, sino el flujo administrativo que circula entre las oficinas y que continúa operando en papel. El sistema propuesto no inscribe, no califica jurídicamente y no certifica; administra, resguarda y da trazabilidad a los documentos que soportan esas funciones.

Los beneficiarios son tres. El personal operativo de cada oficina deja de capturar dos veces los mismos datos y deja de depender del traslado físico para saber si un documento llegó. Los titulares y directivos obtienen, por primera vez, cifras de volumen y de tiempo de atención que pueden verificarse contra el detalle que las originó. Y el ciudadano percibe una reducción en la espera, porque buena parte del tiempo que hoy transcurre no proviene de la calificación jurídica de su asunto sino del desplazamiento y el registro manual del expediente.

Desde el punto de vista técnico, la solución resuelve dos problemas con un mismo mecanismo. La bitácora inmutable de movimientos, exigida para garantizar la trazabilidad, es también la única fuente de datos del módulo de control estadístico. De esta manera, la información directiva se deriva siempre de los hechos registrados y no de capturas paralelas, lo que asegura que el reporte y el detalle coincidan por construcción y no por disciplina del usuario. A ello se suma un criterio de sostenibilidad: resolver el proyecto sobre tecnologías de código abierto y empaquetar el despliegue en archivos de configuración versionados permite que el sistema siga operando al concluir la estadía, mantenido por el área técnica del propio Instituto y sin dependencia del desarrollador. El proyecto no requiere adquisición de licencias, y los ahorros esperados provienen de eliminar la doble captura, reducir el traslado físico y disminuir el tiempo invertido en la integración manual de reportes.

En el plano normativo, el registro auditable de cada acceso y el esquema de permisos diferenciados contribuyen al cumplimiento de las obligaciones que la Ley General de Archivos \cite{lga2018} y la Ley General de Protección de Datos Personales en Posesión de Sujetos Obligados \cite{lgpdppso2017} imponen a la institución. De no realizarse el proyecto, esas obligaciones seguirían atendiéndose sobre un acervo cuyo acceso no deja constancia verificable.

El proyecto es viable ahora porque el Instituto dispone de infraestructura de virtualización propia y porque la solución se construye con herramientas libres, sin costo de licenciamiento. Finalmente, en el plano académico articula las áreas formativas del programa de \ncarrera{} —análisis y modelado de procesos, ingeniería de requerimientos, diseño de bases de datos, desarrollo de aplicaciones web, seguridad informática, administración de servidores y virtualización, y documentación técnica— en un producto que debe operar en un entorno real y ante usuarios reales, que es precisamente lo que el perfil de egreso plantea.



\subsection{Alcances y limitaciones} \label{sec:alcances_limitaciones} % audit:req=alcances

\begin{guia}[Guía: Alcances y limitaciones (Revisión 1)]
\textbf{¿Para qué sirve?} Marca las fronteras del proyecto para evitar expectativas irreales.

\textbf{Alcances} -- lo que \emph{sí} incluye el trabajo. Debe responder: ¿qué módulos, funciones, procesos, usuarios o
datos se cubren? ¿en qué plataformas o entornos? ¿qué entregables habrá? ¿qué queda \emph{explícitamente fuera}?

\textbf{Limitaciones} -- lo que restringe el trabajo y que no depende totalmente de ti. Debe responder: ¿qué límites hay
de tiempo (un cuatrimestre), recursos, acceso a datos o a la infraestructura de la empresa, tecnología, permisos?
¿cómo pueden afectar la validez o generalización de los resultados?

\textbf{Extensión mínima:} 60 palabras en total.

\textbf{Errores frecuentes:} prometer más de lo que cabe en la estadía; confundir limitaciones con excusas; no decir qué
queda fuera.

\textbf{Cómo se revisa:} realismo del alcance y honestidad de las limitaciones.
\end{guia}

\subsubsection{Alcances} \label{subsec:alcances}
% >>> ESCRIBE AQUÍ los alcances <<<

El proyecto comprende el desarrollo completo de una aplicación web con los cuatro módulos funcionales comprometidos en el registro del proyecto: autenticación y control de acceso por rol y oficina de adscripción; solicitud y carga de documentos con validación de formato, tamaño y captura de metadatos; consulta y visualización en línea con bitácora de accesos; y control estadístico con reportes exportables. A ellos se añaden dos componentes de soporte sin los cuales los anteriores no operan: el servicio transversal de bitácora y auditoría, que registra todo evento del sistema y alimenta al módulo estadístico, y la administración de catálogos, usuarios y parámetros, que permite al Instituto operar la solución sin intervención del desarrollador.

Comprende igualmente el diseño e implementación de la base de datos, el esquema de almacenamiento de los archivos fuera de la raíz pública del servidor, y la creación y configuración del entorno virtualizado que aloja la solución, lo que incluye sistema operativo, servicios base, red interna, usuarios, política de respaldo y monitoreo. Los usuarios considerados son el personal de las oficinas del Instituto que participa en el flujo documental interno, bajo seis perfiles: administrador del sistema, responsable de oficina, capturista, consultante, auditor y directivo. El acceso se realiza desde navegador, con interfaz responsiva.

En el plano documental se entregan el diagnóstico del proceso, la especificación de requerimientos, el documento de arquitectura y modelo de datos, los prototipos de interfaz, el plan y las evidencias de pruebas, los manuales técnico y de usuario, y el presente reporte. Se contempla finalmente la capacitación del personal usuario y del área técnica del Instituto.

Queda explícitamente fuera del alcance la digitalización retrospectiva del acervo histórico en papel, por tratarse de un proceso archivístico con requerimientos de equipo y personal ajenos a la estadía; el sistema queda preparado para recibir esa carga posteriormente. También quedan fuera los procesos registrales y catastrales sustantivos, regulados por la legislación aplicable \cite{leyrpp2008, leycatastro2001} y atendidos por el sistema registral electrónico del Instituto, así como la firma electrónica avanzada, el reconocimiento óptico de caracteres y el desarrollo de una aplicación móvil nativa.


\subsubsection{Limitaciones} \label{subsec:limitaciones}
% >>> ESCRIBE AQUÍ las limitaciones <<<

El proyecto se ejecuta durante el periodo de estadía comprendido del 31 de agosto al 4 de diciembre de 2026, con una carga horaria total de 600 horas, lo que acota el alcance a lo que puede construirse, probarse y liberarse dentro de ese lapso y obliga a diferir a trabajo futuro las funcionalidades enumeradas en los alcances.

La limitación de mayor impacto es la dependencia de infraestructura: el proyecto requiere que el Instituto provea el entorno virtualizado con las características mínimas especificadas, y la asignación de ese recurso no depende del estadista sino del área técnica de la institución. Un retraso en su entrega afectaría directamente al cronograma, por lo que se mitiga desarrollando en un entorno local contenedorizado equivalente y migrando después sin cambios en el código.

En segundo lugar, el proyecto opera bajo las políticas de seguridad y de red que el Instituto tenga establecidas, las cuales pueden restringir puertos, accesos remotos o la instalación de componentes, y que no son negociables desde el proyecto.

En tercer lugar, la caracterización del proceso vigente se construyó por análisis documental del flujo y no por medición instrumentada, dado que el proceso actual no genera los registros que permitirían medirlo. En consecuencia, el contraste de resultados en materia de tiempos será cualitativo, mientras que el de trazabilidad sí será verificable a partir de la bitácora del propio sistema. Esta condición limita la generalización de las conclusiones sobre reducción de tiempos.

Finalmente, en ningún caso se emplean datos personales reales en los ambientes de desarrollo y pruebas; la información utilizada es ficticia o proporcionada expresamente por el Instituto para tal fin, lo que restringe las pruebas de volumen a conjuntos sintéticos.

